\input{../tex-inputs/latex-leseansicht-vorspann}

\section[Rainer Maria Rilke an Arthur Schnitzler, 15. 4. 1899]{L04535 Rainer Maria Rilke an Arthur Schnitzler, 15. 4. 1899}
\nopagebreak\mylabel{L04535v}
\rehead{ }\normalsize\beginnumbering\briefempfaengerindex{Schnitzler, Arthur@\textsc{Schnitzler, Arthur}!zzzRilke, Rainer Maria@\emph{von Rainer Maria Rilke}!1899-04-151@{15. 4. 1899}|(be}
\toendnotes[C]{\smallbreak\pagebreak[2]}
\correspDesc{Versand  durch Rainer Maria Rilke am 15. 4. 1899 in Schmargendorf
\newline{}Erhalt  durch Arthur Schnitzler im Zeitraum [16. 4. 1899 – 20. 4. 1899?] in Wien}\toendnotes[C]{\smallbreak}
\Standort{CUL, Schnitzler, B 84.}
\physDesc{Brief, 1 Blatt, 2 Seiten, 839 Zeichen
\newline{}Handschrift: schwarze Tinte, deutsche Kurrent}
\buchAbdrucke{\weitereDrucke{\emph{Rainer Maria Rilke und Arthur Schnitzler. Ihr Briefwechsel.}. Mit Anmerkungen herausgegeben von Heinrich Schnitzler In: \emph{Wort und Wahrheit}, Jg. 13, Nr. 4, April 1958, S. 282–298, hier S. 285.} }\toendnotes[C]{\smallbreak}
\pstart
           {\pb}\textsc{Schmargendorf}{ }\textsuperscript{bei}{ }\textsc{Berlin}\oindex{Schmargendorf@\textbf{Schmargendorf}|pw}\pend
           
\pstart
           \textsc{Villa Waldfrieden\oindex{Villa Waldfrieden@\textbf{Villa Waldfrieden}|pw}}, am 15. April 99.\pend
           \vspace{0.5em}
\pstart
           Daß Sie mir, ſehr verehrter Herr \textsc{Doctor},
               aus verdunkelten Stunden heraus – dennoch{ }ſchreiben, bedeutet mir faſt mehr als ein
               längerer Brief, der{ }ſich leicht und unwillkürlich aus der Stimmung löſt. Vielen
               Dank!\pend
           
\pstart
           Und wenn Sie \label{K_L04535-1v}\edtext{in \textsc{Berlin\oindex{Berlin@\textbf{Berlin}|pw}}}{\lemma{\textnormal{\emph{in Berlin}}}\Cendnote{\textnormal{Schnitzler war zwischen 24. 4. 1899 und 2. 5. 1899 in Berlin\oindex{Berlin@\textbf{Berlin}|pwk}.}}}\label{K_L04535-1}{ }\introOben{}ſind\introOben{}, bitte, erinnern Sie{ }ſich meiner. Vielleicht \label{K_L04535-2v}\edtext{darf ich Sie{ }ſogar im »\textsc{Waldfrieden\oindex{Villa Waldfrieden@\textbf{Villa Waldfrieden}|pw}}« erwarten}{\lemma{\textnormal{\emph{darf … erwarten}}}\Cendnote{\textnormal{Rilke\pwindex{Rilke, Rainer Maria 4.\,12.\,1875 Prag – 29.\,12.\,1926 Montreux@\textsc{Rilke, Rainer Maria} (4.\,12.\,1875 Prag – 29.\,12.\,1926 Montreux)|pwk} reiste zwar erst nach dem
                     25. 4. 1899 nach Russland\oindex{Russland@\textbf{Russland}|pwk} ab, es dürfte aber trotzdem zu keinem Zusammentreffen gekommen
                  sein – zumindest findet sich keines im \emph{Tagebuch}\pwindex{Schnitzler, Arthur 15. 5. 1862 Wien – 21. 10. 1931 ebd.@\textsc{Schnitzler, Arthur} (15. 5. 1862 Wien – 21. 10. 1931 ebd.)!Tagebuch@\strich\emph{Tagebuch}|pwk} erwähnt.}}}\label{K_L04535-2}? – In dieſem beſten Fall beſtimmen Sie auf einer
               Karte die Stunde; jede iſt mir recht.\pend
           
\pstart
           Ich würde nicht wagen bei enger Zeit von \textsc{Schmargendorf\oindex{Schmargendorf@\textbf{Schmargendorf}|pw}} zu reden, wenn \uline{nur}\uline{ich} hier draußen wäre; allein es iſt ja auch der
               erſte Frühling da an meinem Waldrand – und die Gaſſen \textsc{Berlins\oindex{Berlin@\textbf{Berlin}|pw}} wiſſen nichts davon!\pend
           
\pstart
           Am 24. oder 25. reiſe ich fort – aber hoffent{\pb}lich nicht ohne Erinnerung an eine
               Gemeinſamkeit.\pend
           
\pstart
           Von Frau \textsc{Lou Andreas-Salomé\pwindex{Andreas-Salomé, Lou 12.\,2.\,1861 Sankt Petersburg – 5.\,2.\,1937 Göttingen@\textsc{Andreas-Salomé, Lou} (12.\,2.\,1861 Sankt Petersburg – 5.\,2.\,1937 Göttingen)|pw}} Gegengrüße und von mir: Alles Bewuſste!\pend
           
\pstart
           Ihr:{\\[\baselineskip]}\spacefill\mbox{RainerMariaRilke}\pend
           \leftskip=0em{}\selectlanguage{ngerman}\endnumbering\briefempfaengerindex{Schnitzler, Arthur@\textsc{Schnitzler, Arthur}!zzzRilke, Rainer Maria@\emph{von Rainer Maria Rilke}!1899-04-151@{15. 4. 1899}|)be}\mylabel{L04535h}
\begin{anhang}
\end{anhang}\newcommand{\dateiname}{L04535}\newcommand{\titel}{Rainer Maria Rilke an Arthur Schnitzler, 15. 4. 1899}\newcommand{\editorInnen}{Herausgegeben von Jahnke, SelmaMüller, Martin Anton}\input{../tex-inputs/latex-leseansicht-abspann}
